\section{Interpolation and definability}
\label{sec:cut-elimination-interpolation}

Interpolation produces an intermediate formula in the vocabulary shared by
an antecedent and its consequence.  Maehara's partition method and the
relationalization argument below prove it for general first-order signatures.
Beth definability follows by applying interpolation to two expansions of a
common reduct.

\begin{theorem}[Craig interpolation]
\label{thm:craig-interpolation}
Let \(\varphi\) and \(\psi\) be first-order formulas with
\(\entails\varphi\to\psi\).  Then there is a formula \(\theta\) such
that \(\entails\varphi\to\theta\), \(\entails\theta\to\psi\), and
every nonlogical symbol of \(\theta\) occurs in both \(\varphi\) and
\(\psi\); equality and \(\bottom\) may occur in \(\theta\)
regardless.  Moreover,
\[
 \FV(\theta)\subseteq\FV(\varphi)\cap\FV(\psi).
\]
In particular, sentences have a sentence interpolant.
The degenerate cases are covered because \(\bottom\) and
\(\neg\bottom\) are available as interpolants when \(\varphi\) is
unsatisfiable or \(\psi\) is valid.
\textit{Status: proved (Annex~\ref{annex:p7.interp.3}).}
\end{theorem}

The proof uses the partition lemma, direct cut-free completeness for the
relational fragment, and an explicit graph translation for functions and
equality.  Historical treatments are
\cite{Craig1957,Maehara1960,TroelstraSchwichtenberg2000,Hodges1993}.

\begin{theorem}[Beth definability]
\label{thm:beth}
Assume the Axiom of Choice in the ambient set theory, as in
Theorem~\ref{thm:fol-compactness} for arbitrary signatures.
Let \(T\) be a first-order theory in a signature containing a
distinguished relation symbol \(R\) of finite arity \(n\geq0\), and suppose \(T\) defines \(R\)
implicitly: any two models of \(T\) with the same reduct to the signature
without \(R\) interpret \(R\) identically.  Then \(T\) defines \(R\)
explicitly: some formula \(\theta(\bar x)\) not containing \(R\)
satisfies
\[
 T\derives\forall\bar x\,
       (R(\bar x)\leftrightarrow\theta(\bar x)),
 \qquad \FV(\theta)\subseteq\{x_1,\ldots,x_n\}.
\]
The theory and its signature may be infinite, and the signature may contain
functions and constants.  Some finite \(T_0\subseteq T\) already proves
the displayed definition.  For \(n=0\), both the tuple and its quantifier
block are empty.
\textit{Status: proved (Annex~\ref{annex:p7.interp.4}).}
\end{theorem}

The proof is in Annex~\ref{annex:p7.interp.4}; see also
\cite{Beth1953,Hodges1993}.
The model-theoretic route through Robinson joint consistency is proved in
Theorem~\ref{thm:p03-robinson-joint-consistency}, and its equivalence with
Craig interpolation is Theorem~\ref{thm:p03-craig-robinson-beth}.
Beth's theorem complements the conservativity of definitional extensions
(\cref{thm:definitional-conservativity}): adding a defining axiom for
\(R\) produces a theory that defines \(R\) explicitly, hence implicitly,
while proving no new sentences of the smaller signature; Beth definability
states that implicit definability already guarantees an explicit
definition.
